\documentclass[11pt]{article}

\usepackage[margin=1.1in]{geometry}
\usepackage{amsmath,amssymb,amsthm}
\usepackage{booktabs}
\usepackage{microtype}
\usepackage{xcolor}
\usepackage[colorlinks=true,linkcolor=blue!60!black,citecolor=blue!60!black,urlcolor=blue!60!black]{hyperref}
\usepackage{xurl}

\newtheorem{theorem}{Theorem}[section]
\newtheorem{proposition}[theorem]{Proposition}
\newtheorem{corollary}[theorem]{Corollary}
\newtheorem{lemma}[theorem]{Lemma}
\theoremstyle{definition}
\newtheorem{claim}[theorem]{Claim}
\theoremstyle{remark}
\newtheorem{remark}[theorem]{Remark}

\newcommand{\PT}{\mathcal{PT}}
\newcommand{\ii}{\mathrm{i}}
\newcommand{\ee}{\mathrm{e}}
\newcommand{\Tr}{\mathrm{Tr}}
\newcommand{\norm}[1]{\left\lVert #1 \right\rVert}
\newcommand{\ag}{\alpha_g}

\title{Conformal gravity must choose:\\
exceptional-point kinematics, vacuum decay,\\
and the sample space of Weyl gravity}

\author{Jonathan Torkelson\\[2pt]
\small Embeddetech, Inc., Tulsa, Oklahoma, USA\\
\small\texttt{jtorkelson@embeddetech.com}}
\date{July 12, 2026}

\begin{document}
\maketitle

\begin{abstract}
Mannheim's conformal (Weyl) gravity is the flagship application of
the $\PT$/antilinearity program, and its quantum-mechanical prototype
is the Pais--Uhlenbeck oscillator at \emph{equal frequencies}: the
$1/k^{4}$ propagator is, mode by mode, the degenerate limit of two
opposite-residue poles---a Jordan block. We analyze the theory in the
stochastic currency of the preceding papers of this series, with four
results. (i)~\emph{Kinematics}: every linearized spin-2 field mode ($k>0$) of
conformal
gravity
sits exactly at an exceptional point of the similarity dictionary,
whose per-mode cost we lower-bound for every normal-target (mode-built) dictionary (up to an $O(1)$ factor, the EP-driven $\mathrm{cond}\to2/|\delta|$ at $\omega=1$, following the relative
detuning $\delta/\omega$ along a dispersion; massive quadratic gravity approaches the exceptional line
in the ultraviolet with EP-driven cost growing as $k^{2}$), while the
Hermitian ghost realization passes through the point exactly
unitarily. (ii)~\emph{A no-representation lemma}: a real-eigenvalue
Jordan-block
Hamiltonian (finitely many configurations) admits \emph{no} closed stochastic-process representation
under any similarity whatsoever---the obstruction at the conformal
point is structural, not a diverging cost---and the lemma lifts to
the free field level: the $M$-quantum sector of a conformal mode carries a
single Jordan block of size $M{+}1$ (secular growth $t^{M}$, the
constructive form of Smilga's infinite-order exceptional point), and
no metric that is bounded with bounded inverse in the
occupation-number Fock norm yields a stochastic
representation on a conformal mode's Fock
tower. This, we argue, recasts the
published Bender--Mannheim/Smilga dispute: ``unitarity is broken''
is the closed-currency statement, while unitary evolution via
nonstationary states we read, in this currency, as the survival-conditioned
bookkeeping of the dilation---an open-system reading. The dipole
ghost's open process is a one-way cascade whose indivisibility enters, at the
mechanical cell, through conditional mass creation---a mechanism new to the
taxonomy (the higher Fock sectors add ordinary divisor negativity---a
convention-independent sign; the mass-creation reading is stated
relative to the minimal survival-conditioning convention,
whose definition and justification the text makes explicit).
(iii)~\emph{Vacuum decay}: at the scale-invariant ($\delta=0$) point the ghost sector's golden-rule rate sits
on the contact (s-flat) rung of the dimension ladder,
$\Gamma/V\sim\Lambda^{4}/\ag^{2}$---a maximally generous lower bound, since
consistently regulating the split poles ($\delta>0$) only raises it (for modes of relative detuning $\ll1$)---with a dimensionless coupling in
place of the Planck suppression; pricing the Lorentz-violating-cutoff
rescue then forks on $\ag$---the face-value rotation-curve matching
($|\ag|\sim3\times10^{109}$) makes the bound vacuous but that branch fails
solar-system tests on published (Mannheim-disputed) analyses, while
an unfixed $\ag$ within eight decades
of unity (a conservative figure) forces the cutoff to macroscopic wavelengths; the larger-coupling
ray is closed on the face-value branch by point-collapse---every leading-order
gravitational (one-exchange) tree observable scales as $1/\ag$, so the fit pins $\ag$ to a single,
solar-system-dead value---and on Mannheim's branch by an explicit naturalness
prior, stated as such.
(iv)~\emph{The trilemma closes}: coupled ghosts, empirical Lorentz
invariance, definite configurations---pick two (at tree level). Conformal gravity's
no-ghost quantization has already picked: at the exceptional point
its Jordan dynamics admits no \emph{closed}
definite-configurations seat within the stated dictionary classes---a
free/linearized-level statement---so that program is born on
the exit-the-sample-space horn, while the Hermitian ghost
realization keeps configurations at precisely horn 1's price; we
price both seats. The \emph{adjudicated}
escape routes map onto the horns (one horn-adjacent, one premise-disputing; the recent 2024--2026 entries, the older alternative quantizations, and the Bohmian route remain pending).
\end{abstract}

\newpage
\begin{quote}
\small\noindent\textbf{Note on provenance.} This manuscript is an
experiment in AI-conducted research as much as a physics paper. The
project began when the author asked the large language model Claude
Fable~5 (Anthropic): \emph{``Is it not possible that the antilinear
PT/time-reversal symmetry element could maybe predict indivisible
non-Markovian processes?''} Everything that followed was performed by
the model: the research program, the theorems and proofs, the numerical
experiments, the literature search, and this text. The model was
repeatedly prompted to attack its own results adversarially and to
search for prior work. The author, a professional engineer rather than
a physicist, directed and reviewed the sessions but makes no claim to
comprehension of the findings, and is circulating the manuscript for
independent review by professional physicists. Accordingly, nothing
here should be taken on trust: every theorem is stated with a full
proof, and every numerical claim is reproducible from the public
scripts or, where stated, from the in-text analytic computations
(see Code availability); the originating conversation,
project briefing, and full research notebook are preserved in the
program's hub repository,
\url{https://github.com/embeddetech/ghosts-program}.
\end{quote}

\newpage
\tableofcontents

%=====================================================================
\section{Introduction}
\label{sec:intro}

Fourth-order gravity buys renormalizability with a ghost
\cite{Salvio2018}, and every school of quadratic gravity has a
prescription for living with it: Mannheim's conformal program
quantizes the ghost away with $\PT$ methods and zero-norm states
\cite{BenderMannheim2008,BenderMannheim2008PRD,Mannheim2020}; Donoghue and
Menezes keep it and pay in microcausality \cite{DonoghueMenezes};
Anselmi's fakeons remove it by fiat at the level of the propagator
prescription \cite{Anselmi2018}; Salvio and Strumia keep positive
energies at the price of a negative-norm configuration space
\cite{SalvioStrumia2016}; and the agravity analysis of
Gross--Strumia--Teresi--Zirilli, asking whether negative kinetic
energy is even metastable, answers that for gravitationally coupled
ghosts the runaway is ``cosmologically slow''---a Planck-suppression
statement---while conceding that the field-theoretic rate is
``naively infinite'' and that they ``do not know the true rate''
\cite{GSTZ2021}.

The preceding papers of this series developed a currency in which
such prescriptions can be compared \cite{indivisible,instantaneous,conditioned}: quantum theory as indivisible stochastic processes on a
configuration space (Barandes \cite{Barandes2023,Barandes2025}), in
which ghosts are a recurrence-classification problem (conservative
$\to$ transient $\to$ explosive) \cite{indivisible}, the divergent
ghost-vacuum rate of field theory is an instantaneous-state
pathology whose resolution is the choice of sample space
\cite{instantaneous}, and the whole subject obeys a trilemma:
\emph{coupled ghosts, empirical Lorentz invariance, definite
configurations---any two, never all three} (at tree level) \cite{instantaneous}.

This paper aims that currency at the sharpest target: pure conformal
(Weyl) gravity, whose ghost is not a separate massive pole but a
\emph{degeneracy}---the $1/k^{4}$ propagator is, in Mannheim's own
presentation, the limit of two poles of opposite residue as their
masses merge \cite{Mannheim2011}, and the corresponding mechanical
prototype, the equal-frequency Pais--Uhlenbeck (PU) oscillator, is a
Jordan-block Hamiltonian: nondiagonalizable, and---by
Bender and Mannheim's own theorem---possessing \emph{no Hermitian
counterpart} \cite{BenderMannheim2008PRD}. Whether that theory is
unitary is disputed in print: Bender and Mannheim construct unitary
evolution on nonstationary states \cite{BenderMannheim2008PRD};
Smilga concludes that for this nonstandard realization ``unitarity
is broken in the equal frequencies limit'' \cite{Smilga2008}, later
recasting the point as an exceptional point of infinite order
\cite{Smilga2014}; alternative quantizations trade the difficulty
for other costs \cite{Mostafazadeh2011,IlhanKovner2014}; and the
quantization remains contested into this year
\cite{FMT2026,FringTairaTurner2025,FelskiFringTurner2025,DeyFring2026}. Two of the
newest entries sit closest to this paper's machinery and deserve
their coordinates on our map: the ghost-free quantization via
non-unitary similarity transformations~\cite{FringTairaTurner2025}
is a similarity dictionary in precisely the sense of
Sec.~\ref{sec:kinematics}, so at the equal-frequency point it
meets Lemma~\ref{lem:jordan} sector-wise (any sector-preserving Dyson
map), or Theorem~\ref{thm:tower} if its Dyson map is bounded with bounded
inverse in the occupation-number Fock norm (an unbounded, sector-mixing
Dyson map falls outside both, into the open class of
Claim~\ref{claim:choose}); and the Bohmian
analysis of the degenerate cell~\cite{DeyFring2026} equips the
same kinematics with \emph{deterministic} configuration
trajectories---a different beable ontology, whose inequivalence
findings concern trajectory dynamics and are complementary to,
not in tension with, the stochastic-representation lemma proved
here. Beyond the PU quantization dispute, Holdom argues (in $0{+}1$d
models aimed at quadratic gravity) that ghost theories are
intrinsically unitary and perturbatively stable once an appropriate inner product
defines the Born rule~\cite{Holdom2024}---an inner-product
construction this paper's currency would price as a sample-space
exit---the $\PT$ route has been carried to
the full quadratic theory~\cite{Kuntz2024}, a Krein-space
quantization of a four-derivative QFT with a generalized Born rule
has appeared as this manuscript was completed
\cite{BatemanTurok2026}, a direct-sum quantization recasting
the massive tensor ghost as an inverted-harmonic-oscillator
instability held off the asymptotic spectrum has been carried to
the full 4d theory \cite{KumarMarto2026}, the massive ghost
has been argued---conditionally on a ghost--Faddeev--Popov bound
state---to confine into zero-norm BRST-quartet bound
states \cite{Oda2026}, complex-mass ghost pairs have been
argued to bind into states satisfying reflection positivity
\cite{AsoreyKPS2026}, an LSZ scattering construction has been
built atop negative-norm-projecting quantizations of the
quadratic theory \cite{Santillan2025}, and ghost fields have
been argued never to appear as isolated asymptotic particles at
all \cite{Buoninfante2026}; all of these enter the
trilemma registry as pending classifications rather than
adjudicated rows.

Our results, in brief. Section~\ref{sec:kinematics}: the exceptional
point is not a corner of parameter space but the entirety of
conformal gravity's spin-2 kinematics---every mode ($k>0$) of that sector
sits at it, and we
measure what that costs any similarity dictionary with a normal target built on the mode decomposition.
Section~\ref{sec:lemma}: a short lemma shows a real-eigenvalue
Jordan-block
Hamiltonian (finitely many configurations) admits no closed stochastic-process representation under
any similarity; the open (dilated) representation exists and is a
one-way conditioned cascade with a new indivisibility mechanism. The
Bender--Mannheim/Smilga dispute dissolves into two currencies.
Section~\ref{sec:rate}: the vacuum-decay rate of the conformal ghost
sector sits on the contact rung of the series' dimension ladder,
$\Lambda^{4}/\ag^{2}$ (a maximally generous lower bound---regulating the
split poles only raises it for modes of relative detuning $\ll1$), and pricing the Lorentz-violating-cutoff
rescue forks on the value of $\ag$---with neither branch tenable (each on a stated closure---observational death or a naturalness prior).
Section~\ref{sec:trilemma}: the trilemma closes, with every
adjudicated escape route priced against a horn---one route (Merlin)
horn-adjacent, one (metastability) disputing the rate premise rather
than picking a horn---and several recent entries logged as pending.

%=====================================================================
\section{The currency, in one page}
\label{sec:currency}

We use the stochastic--quantum correspondence in the form of
\cite{Barandes2023,Barandes2025,indivisible}: configurations carry the
ontology; dynamics is a family of (generally indivisible) stochastic
maps $\Gamma(t)$. A \emph{closed representation} in a fixed beable
basis is a family $\Gamma(t)=|V(t)|^{2}$ entrywise with
$V(t)=S\ee^{-\ii Ht}S^{-1}$ for a fixed invertible $S$ (the beable
basis is chosen once, so $S$ carries no time dependence), with
$\Gamma(t)$ column-stochastic for all $t$; $\Gamma$ is
\emph{unistochastic} when $V(t)$ is unitary---the case a positive
metric realizes, which is sufficient though not necessary for
stochasticity of $|V|^{2}$. Unbounded or lossy
systems get the \emph{minimal}, generally sub-stochastic
(probability-losing) process, with unitary dilations as the
conservative completions \cite{indivisible}. Three prior results are
load-bearing here. (a)~For diagonalizable $H$ with real spectrum,
the closed representation exists exactly when a positive metric
intertwines $H$ with a Hermitian partner (the converse---that a closed
representation forces such a metric---is \cite{indivisible}'s); its cost is the beable
rotation and metric conditioning \cite{indivisible}. (b)~At an
exceptional point of a $\PT$ dimer the metric collapses with
Petermann scaling while the \emph{open} description passes through
smoothly: the EP is a coordinate singularity of the dictionary, not
of the process \cite{indivisible}. (c)~The trilemma of
\cite{instantaneous}: a coupled ghost sector, empirical Lorentz
invariance, and definite configurations cannot coexist (at tree level); each pair
names a known physics program.

%=====================================================================
\section{Kinematics: conformal gravity's spin-2 sector lives at
the exceptional point}
\label{sec:kinematics}

The quadratic (TT-sector) action of conformal gravity,
$I_{W}=-\ag\int d^{4}x\sqrt{-g}\,C_{\lambda\mu\nu\kappa}
C^{\lambda\mu\nu\kappa}$ with $\ag$ dimensionless, linearizes to
$I_{W}^{(2)}=-\tfrac{\ag}{2}\int(\Box K_{\mu\nu})^{2}$
($K_{\mu\nu}$ the traceless part of the metric perturbation),
giving a
propagator $\propto P_{2}/(\ag k^{4})$ ($P_{2}$ the
transverse-traceless spin-2 projector; the explicit tensor
normalization is our convention, and the overall constant is fixed
below, where the flat-space exchange is written $1/(2\ag k^{4})$;
\cite{Mannheim2011} states the $1/k^{4}$ form). Mannheim presents the $1/k^{4}$ explicitly as
the degenerate limit of two opposite-residue $1/k^{2}$ poles, and
the equal-mass limit as a nondiagonalizable Jordan-block
Hamiltonian \cite{Mannheim2011,Mannheim2020}, whose characteristic
zero-norm states Bender and Mannheim exhibit
explicitly \cite{BenderMannheim2008PRD}.
The linearized particle content---massless spin-2, spin-1, and the
ghost partner---was catalogued by Riegert \cite{Riegert1984}; the
exceptional-point statements of this paper concern the spin-2
sector, where the degeneracy lives.
Mechanically: each spatial mode $k$ is a PU oscillator
\cite{PaisUhlenbeck1950} with frequencies $\omega_{1}(k)=
\omega_{2}(k)=k$---the equal-frequency point, identically in $k$
($k>0$).

What does the degeneracy cost? The first-order mode matrix
$A(\omega_1,\omega_2)$ of the PU oscillator (variables
$q_1=x,\,q_2=\dot x,\,p_1,\,p_2$) has eigenvalues
$\pm\ii\omega_1,\pm\ii\omega_2$ merging at
$\delta=\omega_1-\omega_2=0$ into two $2\times2$ Jordan blocks.
Numerically (\texttt{scripts/t4\_ep\_dictionary.py}): the eigenvector-matrix
condition number---whose divergence is set by the scaling-invariant
angle between the merging eigenvectors, which no column rescaling can
repair, so that up to an $O(1)$ factor it lower-bounds the conditioning of
\emph{any} similarity dictionary with a normal target built on the mode
decomposition (normal modes, isospectral partner, quadratic-level
Bender--Mannheim map---each has conditioning at least the merging pair's
spectral-projector norm, itself $\sim1/\text{angle}$; a normal target
is unitarily diagonalizable and the unitary factor leaves the condition
number unchanged, reducing the bound to this diagonal case; this applies
for $\delta\neq0$, where a normal target exists---at $\delta=0$ there is
none, and Lemma~\ref{lem:jordan} supersedes)---diverges as
\begin{equation}
\mathrm{cond}(V)\;=\;\frac{2.00}{|\delta|}
\qquad(\text{fitted power }{-0.993}\text{ over five decades}),
\label{eq:condscaling}
\end{equation}
(here $V$ denotes the eigenvector matrix of the mode matrix $A$,
distinct from the dictionary evolution $V(t)=S\ee^{-\ii Ht}S^{-1}$ of
Sec.~\ref{sec:currency})
with the far-from-EP control at $\mathrm{cond}=6.8$ and
numerical defectiveness (rank loss to $\sqrt{\text{machine precision}}$,
as expected for a defective pair) at $\delta=0$. The $-1$ power is derived,
not merely fitted: the frequencies are exact eigenvalues of the
family, so the detuning is a \emph{non-generic} unfolding
direction of the exceptional point---the eigenvalue splitting is
linear in $\delta$ and the eigenvector angle within each merging
pair is $\propto\delta$, rather than the generic $\sqrt\delta$ of
a random perturbation. The angle's coefficient is derived, not
only measured: the mode matrix's eigenvector for eigenvalue
$\lambda$ is the chain vector
$(1,\lambda,\beta\omega_1^2\omega_2^2/\lambda,\beta\lambda^2)$
($\beta$ a fixed kinetic normalization),
and linearizing the merging pair of the detuned matrix at
$\beta=\omega=1$ gives eigenvector difference
$\delta\,(0,\ii,\ii,-2)$ against $v=(1,\ii,-\ii,-1)$, hence pair
angle $|\delta|\sqrt{6-1}/2=(\sqrt5/2)\,|\delta|$ (here $6=\|(0,\ii,\ii,-2)\|^2$ and the subtracted $1$ is the squared component along $v$, leaving the perpendicular part; script-emitted:
$\mathrm{angle}/|\delta|=1.11802$ at $\delta=10^{-4}$). The two coefficients are one law at this $\beta=\omega=1$ normalization
(the convention-independent content is the $1/|\delta|$ divergence),
linked by the script-emitted identity
$\mathrm{cond}(V)\times\mathrm{angle}\to\sqrt5$ (computed $2.2363$
against $\sqrt5=2.2361$):
$\mathrm{cond}(V)=\sqrt5/\mathrm{angle}=2/|\delta|$ (the law is
even in the signed detuning). The per-mode
law \eqref{eq:condscaling} is stated at fixed $\omega=1$, where
absolute and relative detuning coincide; along a dispersion the
EP-driven part of the cost follows the dimensionless
\emph{relative} detuning $\delta/\omega$, which is what the
units-only control below isolates. Along the Stelle-type
dispersion of \emph{massive} quadratic gravity
($\omega_1=k$, $\omega_2=\sqrt{k^{2}+m_2^{2}}$), the detuning
magnitude $|\delta(k)|\simeq m_2^{2}/2k$ vanishes in the ultraviolet and the
measured cost grows as $k^{5.0}$, which a units-only control
($\omega_1=2k$: $k^{3.0}$) decomposes into a variable-normalization
factor $k^{3}$ times an \emph{EP-driven excess}
$k^{2.0}$---exactly the inverse relative detuning
$(|\delta|/k)^{-1}\sim(k/m_2)^{2}$. Massive quadratic gravity thus
approaches the exceptional line mode-by-mode in the UV; conformal
gravity, $m_2=0$, sits on it at every $k$.

A one-loop check sharpens the scope of this statement. The ghost
pole acquires a width---$\mu^{2}=m_r^{2}+\ii\gamma$ with
$\gamma=\xi^{2}m_r^{2}N_{\text{eff}}/(320\pi)$ ($m_r$ the real pole
mass, $\xi$ the ghost--matter coupling, $N_{\text{eff}}$ the
effective loop multiplicity)~\cite{DonoghueMenezes}---and one might
expect the width to detune
the theory off the exceptional line. It does not: the shift is
Lorentz-invariant in $q^{2}$, so at spatial momentum $k$ the
frequency-space width $\gamma/2E_k$ time-dilates away at exactly
the $1/k$ rate of the real detuning magnitude $m_2^{2}/2k$; both
components of the complex detuning, of magnitude
$|m_2^{2}+\ii\gamma|/2k$, vanish together,
and the cost still diverges, suppressed only by the $k$-independent
factor $[1+(\gamma/m_2^{2})^{2}]^{-1/2}$
($\approx1-2\times10^{-4}$ at $\xi=1$,
$N_{\text{eff}}\approx20$, with $m_r\simeq m_2$ near the exceptional line). What the loop does change is the
approach: the detuning ray is tilted into the complex plane by the
fixed angle $\arctan(\gamma/m_2^{2})$, and the coalescing pair is
(real, complex) at every finite $k$---a resonance-type exceptional
point with unequal linewidths, not a real-axis $\PT$-symmetric one.
We flag explicitly that the saturation a fixed-linewidth
(cavity-photonics) intuition would predict,
$\mathrm{cond}_{\max}\sim2\omega/\Gamma_\mu=640\pi/(\xi^{2}
N_{\text{eff}})$ (with $\Gamma_\mu\equiv\gamma/m_r$ the frozen pole
width as a frequency), is an artifact of freezing the width in
frequency; the relativistic dispersion removes it. Scope: this is
the one-loop, weak-coupling ($m_2\ll M_{\rm Pl}$) pole statement; off-shell
mode matrices at $|q^{2}|\gg m_2^{2}$ remain tree-level claims.

Meanwhile the Hermitian ghost realization---the
ghost-spectrum theory on real configurations $(x,\dot x)$, a
\emph{different} operator not similar to the confluent Jordan
one---passes
through the exceptional point exactly: on a grid at
$\omega_1=\omega_2$, $\norm{H-H^{\dagger}}=8\times10^{-15}$ and
$\Gamma=|U|^{2}$ is doubly stochastic to $1.1\times10^{-14}$
(re-verifying \cite{indivisible} at the field-kinematics framing).
The singularity is the dictionary's alone. The dictionary, however,
is what the no-ghost program \emph{is}; the next section shows the
failure at $\delta=0$ is not a divergence but a nonexistence.

%=====================================================================
\section{No closed sample space at the conformal point}
\label{sec:lemma}

\begin{lemma}[no closed process at a Jordan point]
\label{lem:jordan}
Let $H$ be a matrix on finitely many
configurations with a Jordan block of size $m\ge2$ at a
real eigenvalue (hence nondiagonalizable). Then there is no invertible $S$ for which
$\Gamma(t)=|S\ee^{-\ii Ht}S^{-1}|^{2}$ is column-stochastic---in
particular, unistochastic or doubly stochastic---for all $t$: no
metric, no partner, no beable rotation yields a closed stochastic
representation of any kind.
\end{lemma}

\begin{proof}
If $\Gamma(t)=|V(t)|^{2}$ is column-stochastic on $n$
configurations, then
\[
\sum_{ij}|V_{ij}(t)|^{2}
=\sum_{j}\bigl(\text{$j$th column sum of }\Gamma\bigr)
=\sum_{j}1=n,
\]
so $\norm{V(t)}_{2}\le\norm{V(t)}_{F}=\sqrt n$ is bounded in $t$. But
$\norm{\ee^{-\ii Ht}}\gtrsim t^{\,m-1}$ is unbounded (real eigenvalue
$\Rightarrow|\ee^{-\ii\lambda t}|=1$, so the Jordan
nilpotent part is undamped), and
$\norm{V(t)}\ge\norm{\ee^{-\ii Ht}}/\mathrm{cond}(S)$ is unbounded
for every invertible $S$.
\end{proof}

\begin{remark}
The correspondence itself requires only per-time unistochasticity
\cite{indivisible}---a weaker property than $V(t)$ being unitary,
since unistochasticity fixes only the moduli of some unitary, not
of $V$. The lemma is therefore stated for the weakest relevant
condition (column-stochasticity of $|V|^{2}$) and excludes the
unistochastic representations a fortiori. The core mechanism---column-stochasticity
pinning $\norm{V(t)}_{F}=\sqrt n$ against the unbounded Jordan growth
of $\norm{\ee^{-\ii Ht}}$ under any fixed similarity---appears for the
size-2 boundary case ($s=s_*$) in \cite{indivisible}; the lemma extends
that exclusion to arbitrary block size $m$, to any real eigenvalue,
and, below (Theorem~\ref{thm:tower}), to the Fock tower.
\end{remark}

Numerically (\texttt{scripts/t4\_jordan\_process.py}): growth exponents
$1.000,\,1.999,\,2.997$ for $m=2,3,4$ (with prefactors converging
to $1/(m{-}1)!$); the similarity inequality
verified with margin; the detuned control saturates at
$\norm{U}\sim2/\delta$ as diagonalizability requires. (A numerical
remark worth recording: machine-epsilon perturbations split a
Jordan block by $\sim\epsilon^{1/m}$ and convert polynomial to
exponential growth for $t\gtrsim\epsilon^{-1/m}$; the fragility is
itself the exceptional-point physics.)

Lemma~\ref{lem:jordan} upgrades the kinematics of
Sec.~\ref{sec:kinematics}: away from the conformal point the
dictionary exists at diverging cost \eqref{eq:condscaling}; at it,
the dictionary does not exist. Since Bender and Mannheim themselves
established that the equal-frequency Hamiltonian ``has no Hermitian
counterpart'' \cite{BenderMannheim2008PRD}---their diagnosis of
exactly the nondiagonalizability on which the lemma bites---the
lemma's content in
context is the missing corollary: \emph{for the Jordan dynamics,
whose lack of a Hermitian counterpart is its nondiagonalizability,
there is no closed probability representation on any beables}---under
any similarity at the (finite-dimensional) quadratic level, sector by
sector on the tower, and, by the tower theorem below, under any
bounded dictionary with bounded inverse on
the Fock space. The sector-level exclusion carries no boundedness
hypothesis---each $M$-quantum sector, like the mechanical mode matrix
itself, is finite-dimensional---so it already covers the
quadratic-level avatars of the realized dictionaries and every
sector-preserving field dictionary; the
bounded-inverse hypothesis of Theorem~\ref{thm:tower} buys the
further generality of bounded field dictionaries that mix sectors. For the
no-ghost quantization, the option the
trilemma calls ``definite
configurations'' is thus, within the stated dictionary classes,
structurally absent at the conformal point, not merely
expensive---while the Hermitian ghost realization of
Sec.~\ref{sec:kinematics}, a \emph{different} operator not similar
to the Jordan one, keeps definite configurations together with the
ghost spectrum, which is exactly the seat
Sec.~\ref{sec:rate} prices.

\subsection{From the mechanical prototype to the Fock tower}
\label{sec:tower}

The lemma so far concerns one mode. Second quantization makes it a
field statement, constructively.

\begin{proposition}[sector tower]
\label{prop:tower}
Let the one-particle matrix of a mode be a size-2 Jordan pair,
$A=\omega\mathbf 1+N$ with $\omega\in\mathbb R$, $Nv=u$, $Nu=0$
(the equal-frequency PU
cell). Then $H=d\Gamma(A)$ preserves each $M$-quantum sector
(dimension $M{+}1$), and on it carries a \emph{single} Jordan block
of size $M{+}1$ at eigenvalue $M\omega$: secular growth $t^{M}$,
with degree unbounded up the tower.
\end{proposition}

\begin{proof}
$d\Gamma(N)$ acts on the occupation basis
$\mathrm{sym}(u^{\nu}v^{M-\nu})$ as a weighted shift $\nu\to \nu+1$ with
nonzero weights, so $\bigl(d\Gamma(N)\bigr)^{M}$ maps
$\mathrm{sym}(v^{M})$ to a nonzero multiple of
$\mathrm{sym}(u^{M})$ while $\bigl(d\Gamma(N)\bigr)^{M+1}=0$; the
kernel of $d\Gamma(N)$ in the sector is one-dimensional (geometric
multiplicity 1).
\end{proof}
Verified numerically for $M\le8$: nilpotent ranks,
geometric multiplicities, and fitted growth exponents
$1.000,2.000,3.999,5.999,7.999$ at $M=1,2,4,6,8$
(\texttt{scripts/t4\_sector\_tower.py}).

\begin{remark}[the sector similarity is positive-diagonal]
\label{rem:posdiag}
The shift of Proposition~\ref{prop:tower} is not merely nonzero- but
\emph{positively}-weighted: on the occupation basis $\mathrm{sym}(u^{\nu}v^{M-\nu})$
(taking $\{u,v\}$ orthonormal---the convention defining the occupation beables),
$d\Gamma(N)=a_{u}^{\dagger}a_{v}$ sends
$|\nu\rangle\mapsto\sqrt{(\nu{+}1)(M{-}\nu)}\,|\nu{+}1\rangle$, with strictly positive
weights for $\nu=0,\dots,M{-}1$. Hence $d\Gamma(N)=D\,J_{M}\,D^{-1}$ with $J_{M}$
the standard size-$(M{+}1)$ nilpotent block (ones on the subdiagonal, matching
the shift $\nu\to \nu{+}1$) and $D=\mathrm{diag}(d_{\nu})$, $d_{0}=1$,
$d_{\nu+1}=\sqrt{(\nu{+}1)(M{-}\nu)}\,d_{\nu}$ ($\nu=0,\dots,M{-}1$), a \emph{fixed,
positive} diagonal (\texttt{scripts/t4\_tower\_similarity.py}). The occupation
and Jordan bases therefore differ, at every $M$, by one $t$-\emph{independent}
positive-diagonal similarity: $U(t)=e^{-\ii M\omega t}\,D\,e^{-\ii J_{M}t}D^{-1}$,
so $|U(t)|^{2}=D^{2}\,|e^{-\ii J_{M}t}|^{2}\,D^{-2}$ entrywise (the global phase
$e^{-\ii M\omega t}$ cancels), and any two-time object $\Gamma(T)\Gamma(s)^{-1}$ built from
$\Gamma=|U|^{2}$ conjugates by the same $D^{2}$, scaling entry $(i,j)$ by
$(d_{i}/d_{j})^{2}>0$: at each $M$ the beable-basis divisor shares the entrywise
sign pattern of the Jordan-basis one. Moreover $e^{-\ii J_{M}t}$ is triangular
\emph{Toeplitz}---its offset-$\ell$ entry is $(-\ii t)^{\ell}/\ell!$, independent of
$M$---so the Jordan-basis divisor's offset-$\ell$ entry is one and the same at every
$M\ge \ell$; a depth-$\ell$ sign shown for one block size therefore holds, in the
beable basis, at every sector $M\ge \ell$. ($t$-independence is essential: it is what lets
$D^{2}$ pass through the two-time divisor $\Gamma(T)\Gamma(s)^{-1}$.)
\end{remark}

The hypothesis is Sec.~\ref{sec:kinematics}'s kinematics: $A$ is
the one-particle \emph{Hamiltonian} of the equal-frequency pair
(real eigenvalue $\omega$), which inherits the nilpotent structure
$N$ of the mode matrix's $+\ii\omega$ Jordan cell (the
positive-frequency one of its two
conjugate $2\times2$ blocks---the eigenvalue is the cell's times
$\ii^{-1}$, the nilpotent the same up to that scalar, the Jordan
structure unchanged; this per-level Jordan structure is the
Bender--Mannheim equal-frequency construction \cite{BenderMannheim2008PRD}),
so the Fock tower of
every conformal mode carries this structure.
This is the constructive form of Smilga's observation that the
equal-frequency point is an exceptional point of \emph{infinite
order}, whose closure produces a continuous spectrum
\cite{Smilga2014}. Two consequences. First, no truncation escapes
the lemma---every level $M\ge1$ of the tower contains Jordan
structure, and
each level makes the secular degree worse. Second, the
no-representation statement lifts to the field:

\begin{theorem}[no boundedly invertible dictionary on the Fock tower]
\label{thm:tower}
Let $H=d\Gamma(A)$ as above ($A=\omega\mathbf 1+N$ a size-2 Jordan
pair with $\omega\in\mathbb R$) on the Fock space of one mode
(sector-preserving with finite-dimensional sectors, so
$\ee^{-\ii Ht}$ is defined sector-wise on a dense domain; it is
\emph{not} bounded on $\ell^{2}$---its sector norms grow
exponentially in $M$ at fixed $t$), and
let $S$ be any bounded operator with bounded inverse on the
$\ell^{2}$ configuration space. Then
$\Gamma_{ij}(t)=|(S\ee^{-\ii Ht}S^{-1})_{ij}|^{2}$ cannot have unit
column sums for all $t$ (whenever the representation exists at all,
i.e.\ $\ee^{-\ii Ht}S^{-1}e_{j}\in\ell^{2}$, which we take as part of the
dictionary class; an entrywise column that leaves $\ell^{2}$ is not a
bounded-in-Fock-norm representation and is excluded by definition).
\end{theorem}

\begin{proof}
Unit column sums mean $\norm{V(t)e_{j}}=1$ for every basis vector
$e_{j}$ and all $t$. Let $u,v$ be
the Jordan pair in the one-quantum sector, and let $u^{*}$ and
$\alpha$ be the coordinate functionals dual to $u$ and $v$ there,
extended by zero off the sector (bounded); since
$\ee^{-\ii Ht}$ preserves sectors, they read off the one-quantum
component, on which $H$ acts as $A$, giving
$\langle u^{*},\ee^{-\ii Ht}\psi\rangle
=\ee^{-\ii\omega t}\bigl(\langle u^{*},\psi\rangle
-\ii t\,\alpha(\psi)\bigr)$
for every $\psi$ in the domain.
If $\alpha(S^{-1}e_{j})=0$ for every $j$, then the bounded
functional $\alpha\circ S^{-1}$ vanishes on a basis, hence
everywhere, hence $\alpha=0$ (as $S^{-1}$ is surjective)---false.
So some $j$ has $\alpha(S^{-1}e_{j})=a\neq0$, and with
$c=|\langle u^{*},S^{-1}e_{j}\rangle|$,
\[
1=\norm{V(t)e_{j}}
\ge\frac{\norm{\ee^{-\ii Ht}S^{-1}e_{j}}}{\norm{S^{-1}}}
\ge\frac{|a|\,t-c}{\norm{u^{*}}\,\norm{S^{-1}}}
\;\longrightarrow\;\infty .
\qedhere
\]
\end{proof}

\begin{remark}[mode-wise extension to the full free Fock space]
\label{rem:modewise}
The theorem extends to the multi-mode free field:
$H=\sum_{k}d\Gamma(A_{k})$ commutes with every per-mode number
operator $\hat N_{k}$, so the two-dimensional subspace \{one quantum
in mode $k$, vacuum elsewhere\} is invariant and lies in an invariant
orthogonal decomposition; the dual functionals extend by zero and the
linear-growth argument of Theorem~\ref{thm:tower} runs verbatim
mode-wise.
\end{remark}

A control makes the mechanism sharp: for a \emph{diagonalizable}
one-particle matrix at the same degenerate frequency, every sector
evolution is exactly unitary ($\bigl|\norm{\ee^{-\ii d\Gamma
t}}-1\bigr|$ at machine epsilon, $2.2\times10^{-16}$, at all
sectors and times tested):
degeneracy alone is harmless; the Jordan pair is what kills the
tower. Scope, stated plainly: this is the \emph{free} field on flat
space---precisely the arena whose per-mode kinematics
Sec.~\ref{sec:kinematics} measured; interactions and curved
backgrounds remain outside the theorem. ``Bounded'' here is relative to
the occupation-number ($\{u,v\}$-orthonormal) Fock norm of
Remark~\ref{rem:posdiag}: a Dyson map bounded in a different Fock norm
falls into the honestly-open unbounded/sector-mixing class of
Claim~\ref{claim:choose}, not around the theorem.

\subsection{What survives: the open process, and the adjudication}

What the indivisible-process paper~\cite{indivisible} found at the dimer exceptional point holds here: the
\emph{open} route is untouched (for the mechanical cell, and sector by
sector on the tower---$H$ preserves each finite sector, each carrying
its own contraction and dilation; no tower-wide normalization exists,
by the unbounded sector norms of Theorem~\ref{thm:tower}). Normalizing
$C(t)=\ee^{-\ii Ht}/\sigma_{\max}(t)$, with $\sigma_{\max}$ the
largest singular value, gives an exact contraction;
its Halmos dilation is unitary (verified to $1.4\times10^{-8}$,
limited by a
matrix square root at an exactly-degenerate singular value); and the
reduced process is sub-stochastic at all times. The normalization is
a convention, and the minimal one: a unitary dilation exists exactly
when the normalized evolution is a contraction, so $\sigma_{\max}$
is the pointwise-least normalization admitting a conservative
completion---the unique choice that discards no survival probability
beyond what the dynamics itself loses ($\norm{C(t)}=1$: the
least-decaying state is unconditioned at every time). Its structure is,
however, new to the series' taxonomy: for the pure Jordan cell,
$|U(t)|^{2}$ is \emph{triangular}---the beable-level process is a
one-way cascade into the growth direction---and its indivisibility
enters not through divisor negativity (the broken-dimer mechanism of
\cite{indivisible}; the dimer at its own exceptional point develops a
negative divisor entry by $t\approx0.7$ at unit coupling ($\kappa$, the
cell's off-diagonal coupling)---the EP
limit $1/(\kappa\sqrt2)$ of the companion's onset formula; the
companion's canonical $\theta=\pi/4$ device gives
$t\approx1.0$---while the size-2 PU cell never
does, its $2{\times}2$ divisor being too small to carry a depth-two entry) but through the
divisor violating the column-sum
condition: \emph{conditional mass creation}, onset at finite length
($t_c=0.486$ for the unit cell in the zero-step limit---the root
of $2t\,\sigma_{\max}^2(t)=(\sigma_{\max}^2)'(t)$, emitted by the
script;
$\approx0.50$ at this section's protocol, whose base grid and pair
step are both $0.05$), with
entrywise nonnegativity
intact throughout \emph{for the size-2 cell} (whose triangular divisor has
nonnegative entries). At the higher Fock sectors $M\ge2$ the divisor's
depth-two entry $\tfrac14(T^{2}-s^{2})(T^{2}-3s^{2})$ turns negative for a short
step ($s<T<\sqrt3\,s$), carried into the beable basis by the positive-diagonal
similarity of Remark~\ref{rem:posdiag}---so there the companion's ordinary
negativity mechanism reappears alongside mass creation. One bookkeeping fact must be
flagged, because the two mechanisms differ in kind under it:
divisor entry \emph{signs} are invariant under every positive
scalar renormalization $c(t)$ of the normalization, while divisor
\emph{column sums} scale by $c^{2}(s)/c^{2}(s+\tau)$---negativity
is convention-free, but the mass-creation onset $t_c$ is stated
relative to the minimal ($\sigma_{\max}$) convention. The
convention is not idle, and its minimality carries the physics: for
the unit cell with coupling $\kappa$ the raw $|U|^{2}$ divisor's
column sum is exactly $1+\kappa^{2}(T^{2}-s^{2})$, so a
renormalization avoids the crossing at all step pairs only if
$c^{2}(t)\ge c^{2}(0)\,\ee^{\kappa^{2}t^{2}}$, and that Gaussian
envelope does avoid it (both verified in
\texttt{scripts/t4\_jordan\_process.py}, block~[3b]). Divisibility can thus
be purchased, but only by conditioning away survival probability at
a Gaussian rate against polynomially growing dynamics---a
conditioning overhead $c^{2}/\sigma_{\max}^{2}$ that diverges;
under the minimal convention, in which conditioning means survival
and nothing else, the crossing is unavoidable. The dipole ghost,
read as a process, is a one-way
conditioned cascade.

This dissolves the published dispute into two currencies:

\begin{itemize}
\item \emph{Smilga's statement} (said of the nonstandard
realization: ``unitarity is broken in the equal frequencies limit''
\cite{Smilga2008}) is the closed-currency fact,
now a lemma: no similarity restores a unitary---hence
unistochastic---closed evolution.
\item \emph{Bender--Mannheim's construction} (unitary time evolution
carried by nonstationary states \cite{BenderMannheim2008PRD}) is,
in this currency, the conditional bookkeeping of the dilated open
process (read sector-wise on the tower): the polynomial norm growth is
divided out exactly as the
broken-phase dimer's exponential growth was in \cite{indivisible}---a
survival-conditioned, open-system reading, not a closed sample
space.
\end{itemize}

Both are right where they compute; neither headline survives the
translation. We anticipate the natural reply from the $\PT$ side:
in the theory's own inner product nothing is conditioned, and
probability is exactly conserved. Agreed---and that is precisely
the point at issue: the conditioning is relative to the
Dirac/beable norm, the only norm that assigns Born weights to
\emph{configurations}, so the disagreement is not about the
formalism but about whether a configuration sample space exists.
That is the trilemma's third seat, not a technicality. And the dictionary line for field theory: the Jordan
$t$-term in amplitudes---the $t^{2}$ in probabilities---is precisely
the notorious per-mode linear growth of fourth-order-gravity
perturbations \cite{HellLustZoupanos2023}: conformal
gravity is the
everywhere-realized member of the exceptional-point polynomial class
of \cite{indivisible}.

%=====================================================================
\section{Vacuum decay: the contact rung and the $\ag$ fork}
\label{sec:rate}

Suppose one refuses the exit and keeps definite configurations: the
ghost sector is then subject to the vacuum-decay analysis of
\cite{instantaneous}. There the golden-rule rate factorizes into an
infinite boost-orbit volume times an invariant reduced integral, and
the reduced measure's ultraviolet scaling is set by the vertex:
contact $|\mathcal M|^{2}\sim1$ gives $\Lambda^{4}$, the
graviton-mediated phantom of Cline--Jeon--Moore,
$|\mathcal M|^{2}\sim s^{2}/M_{\rm Pl}^{4}$, gives
$\Lambda^{8}/M_{\rm Pl}^{4}$ \cite{CJM2004,instantaneous}.

For conformal gravity the physical channel is the
Cline--Jeon--Moore analog, vacuum $\to$ (ghost-graviton pair) $+$
(photon pair), through the $1/(\ag k^{4})$-scaled propagator; the
matter side couples through the standard linearized
$\tfrac12 h_{\mu\nu}T^{\mu\nu}$ coupling (the quadratic-action form
of the sourced field equation of~\cite{Mannheim2011}). The
illustrative $T$-parallel, with pair matrix elements
$\langle\text{pair}|T|0\rangle\sim s$,
\begin{equation}
\mathcal M\;\sim\;s\cdot\frac{1}{\ag s^{2}}\cdot s\;\sim\;
\frac{1}{\ag},
\qquad
|\mathcal M|^{2}\sim\frac{1}{\ag^{2}}\;(\text{$s$-flat}),
\end{equation}
carries the correct $\ag$ counting for the physical channel as well
(the same $\ag^{-1}$ follows for the ghost-self-interaction route by
$\ag$ bookkeeping: each graviton vertex drawn from the $\ag C^{2}$
action carries $\ag$, each internal line $1/\ag$, and each
canonically normalized external leg $\ag^{-1/2}$, so a connected
$0\to4$ tree with $n_v$ vertices and $n_v{-}1$ internal lines scales as
$\ag^{\,n_v-(n_v-1)-2}=\ag^{-1}$, independent of topology), and the
$s$-flatness is \emph{forced} by scale invariance: pure conformal
gravity has no dimensionful parameter, so any tree-level $0\to4$
massless amplitude is $s^{0}f(\ag;t/s,u/s)$---scale invariance fixes
the overall Mandelstam degree, not the angular dependence (assumed
integrable; any collinear/soft non-integrability would only raise the
rate, consistent with the lower-bound reading). The reduced measure therefore
sits on the \emph{contact} rung at the scale-invariant point---with
$\Gamma$ now the vacuum-decay rate (not the Sec.~2 stochastic map) per
volume $V$:
$\Gamma/V\sim c_{\rm ps}\,\Lambda^{4}/\ag^{2}$---a maximally generous
lower bound, since consistently regulating the split poles only raises
it (for modes of relative detuning $\ll1$)---with
$c_{\rm ps}=(2\pi)^{-4}$ the phase-space prefactor used throughout
(slopes validated against \cite{instantaneous} to $10^{-3}$;
\texttt{scripts/t4\_rate\_ladder.py}). The Planck \emph{mass} has left the
suppression: what remains is whatever $\ag$ is. This is the
structural sense in which the agravity escape---``cosmologically
slow \dots\ where ghosts have gravitational interactions only''
\cite{GSTZ2021}---does not port to conformal gravity: that estimate
is a statement about $M_{\rm Pl}$, and conformal gravity has none in
its ghost vertex.

Pricing the Lorentz-violating-cutoff rescue (horn 1 of the
trilemma) then requires a value of $\ag$, and here the theory
forks, because no numerical $\ag$ exists in print:

\emph{Face-value branch.} Matching the flat-space
$1/(2\ag k^{4})$ exchange to the Mannheim--O'Brien rotation-curve
parameters \cite{MannheimOBrien2011} gives $|\ag|\sim3\times10^{109}$
(two normalization conventions differing by a fixed factor $4/3$; the in-print sign
constraint is negative \cite{Mannheim2015}). At such $\ag$ the
energy-density bound on the cutoff is vacuous
($\Lambda_{\max}\sim5\times10^{7}\,M_{\rm Pl}$). But this same
face-value
flat-space theory is the one shown to fail solar-system tests
outright in the weak-field limit (wrong light bending \cite{Flanagan2006}, disputed by
Mannheim \cite{Mannheim2007}; wrong Newtonian
limit \cite{Yoon2013}, disputed in \cite{Mannheim2015}): \emph{the
branch that survives the vacuum-decay bound as stated (the contact-rung
minimum) is already dead by other observations---or, if Mannheim's
rebuttals stand (they invoke the very conformal-source structure that defines
the Mannheim branch), one is thereby on the other branch.}

\emph{Mannheim branch.} Mannheim rejects the face-value inversion
(conformal sources; the local Higgs field $S(x)$)
\cite{Mannheim2015}; then $\ag$ is unfixed, and for any $\ag$
within eight decades of unity (a conservative figure; the same bound in fact
reaches $|\ag|\lesssim3\times10^{12}$, below) the same energy-density bound
(method anchored on the phantom vertex, where it reproduces the
published gamma-ray bound scale within an order of magnitude:
$108$ vs $18.8$~MeV, the latter being the $95\%$~C.L.\ limit
quoted in \cite{ClinePhantom2024}; using the critical
density $\rho_{\rm cr}$ rather than the ${\sim}10^{4}\times$
smaller radiation background is the \emph{weakest defensible}
bound, i.e.\ maximally generous to the theory) forces
\begin{equation}
\Lambda_{\max}=\Bigl(\frac{\ag^{2}\rho_{\rm cr}}
{t_{0}\,c_{\rm ps}}\Bigr)^{1/5}
\;\sim\;10^{-8}\,\text{eV}\;\;(|\ag|=1),
\end{equation}
with $t_{0}$ the age of the universe:
a ghost sector cut off at a reduced wavelength
$\hbar c/\Lambda_{\max}\approx20\,$m (physical wavelength
${\approx}130\,$m)---no laboratory-scale field content at all.

\emph{The gap between the branches, priced.} The two branches
bracket the positions actually taken in print; they do not exhaust
the real line, and we close the gap explicitly rather than by
silence. Since the bound scales as
$\Lambda_{\max}\propto\ag^{2/5}$, laboratory-scale field content
($\Lambda_{\max}\gtrsim\mathrm{meV}$) returns for
$|\ag|\gtrsim3\times10^{12}$, and a particle-scale cutoff
($\Lambda_{\max}\sim100$~MeV) at $|\ag|\sim10^{40}$: couplings with
$|\ag|\gtrsim3\times10^{12}$---the entire unanchored ray above threshold,
bounded by nothing from above---would evade the energy bound as stated.
That ray is nonetheless closed, and differently on the two branches. On the
\emph{face-value} branch the closure is observational and needs no prior:
every leading-order (one-exchange) gravitational observable carries a single
$1/(2\ag k^{4})$ propagator and so scales as $1/\ag$ (nonlinear observables
scale as higher powers of $1/\ag$ and close even faster), whence the linear
potential coefficient
$\gamma_*\propto1/\ag$ is \emph{measured} by the rotation-curve fit to be
$|\ag|\simeq3\times10^{109}$---itself already dead by solar-system tests
\cite{Flanagan2006} (a dark-matter-independent, though Mannheim-disputed, anchor)---and any $|\ag|$ off
that value fails galactic and
solar dynamics monotonically (below it the linear term overshoots the
observed pull, above it undershoots by the matching decades). On the
\emph{Mannheim} branch $\ag$ decouples---on Mannheim's contention, via
degeneracy with the undetermined conformal-source normalization---from every conformally-invariant
\emph{classical} observable, so no classical or astrophysical-fit measurement
constrains it (graviton--graviton scattering does depend on $\ag$, but is
unobservable here); there the ray is closed by a
\emph{naturalness prior}, stated rather than smuggled: a dimensionless
coupling is not taken tens of decades above unity without a mechanism. (The
Mannheim energy bound already reaches the threshold on its own---%
$\Lambda_{\max}\lesssim\mathrm{meV}$ for all $|\ag|\lesssim3\times10^{12}$---%
leaving only the ray above it to the prior; and replacing $\rho_{\rm cr}$ by
the smaller measured radiation backgrounds tightens $\Lambda_{\max}$ by only
the fifth root of the density ratio, a factor of order ten.) Under either
branch's closure horn~1 is untenable (on the Mannheim branch up to the stated
prior); absent any closure it survives only on
the unanchored ray $|\ag|\gtrsim3\times10^{12}$, selected to evade the
bound.

Neither branch keeps horn 1 alive, and the unanchored ray of larger
couplings is closed on each (observationally on the face-value branch, by
the stated prior on Mannheim's). And the escape Mannheim himself
offers from both---the $\PT$/zero-norm quantization, under which
the decay channel closes identically
\cite{BenderMannheim2008,Mannheim2020}---is not a fourth option:
it is horn 3, and Sec.~\ref{sec:lemma} has already priced it.

\emph{Scope.} Everything in this section is tree-level, golden-rule,
and computed on the Minkowski background with Dirac-norm (Fock)
quantization, with the exactly degenerate external pair understood as
the smooth $\delta\to0$ limit of the split-pole spectrum---an assumption,
since the merging poles carry opposite residues (the
degenerate tower of Prop.~\ref{prop:tower} has no stationary
final-state basis at $\delta=0$)---that is, on the multiplicity (instantaneous-decay-products) sample-space choice of
\cite{instantaneous}, which is exactly the reading in which horn 1 is
being priced. Indeed the assumption is not merely unproven but fails in
the regulated theory: canonically normalizing the split ($\delta>0$)
external legs against their magnitude-$1/(\ag\delta)$ residues multiplies
each gravitational leg by a divergent factor $\propto\delta^{-1/2}$, so
the consistently regulated golden-rule rate is parametrically larger than the
contact-rung value (for modes of relative detuning $\ll1$) and has no finite $\delta\to0$ limit---the split-pole
and $\delta=0$ computations sit on opposite sides of a non-commuting
limit. Because these Dirac-norm channel rates are positive and cannot
cancel (positivity is norm-relative), this only tightens (never loosens) the horn-1 pricing; the contact rung is
the maximally generous (minimum) reading. It is not a statement about the $\PT$-quantized
theory, whose channels close by construction; that is the point of
the fork, not an oversight of it.

%=====================================================================
\section{The trilemma, closed}
\label{sec:trilemma}

\begin{claim}[conformal gravity must choose---and already has]
\label{claim:choose}
Of the trilemma's three virtues---a coupled ghost sector, empirical
Lorentz invariance, definite configurations---conformal gravity's
\emph{no-ghost} kinematics removes the third, within the stated
dictionary classes, at the outset:
by Lemma~\ref{lem:jordan} the equal-frequency (Jordan) dynamics
admits
no closed configuration-space representation under any dictionary in the
stated classes---any similarity per sector, or any dictionary bounded with
bounded inverse in the occupation-number Fock norm (Theorem~\ref{thm:tower}, extended mode-wise to the full free Fock space by Remark~\ref{rem:modewise}). (Unbounded,
sector-mixing field dictionaries, continuum-beable $(x,\dot x)$ kernels on
$L^{2}$, and the deterministic Bohmian configuration trajectories
of~\cite{DeyFring2026} fall outside these classes and remain open.) Within the stated classes, definite
configurations survive only in the Hermitian ghost realization,
never jointly with the other two virtues (at tree level).
The remaining choices are priced:
\begin{itemize}
\item \textbf{Horn 1} (coupled ghosts $+$ configurations $\Rightarrow$
  Lorentz-violating frame): the $\ag$ fork of
  Sec.~\ref{sec:rate} (priced under the smooth $\delta\to0$ final-state
  limit of that section's Scope)---vacuous-but-dead, or a
  macroscopic-wavelength cutoff. Untenable on either branch (Mannheim up to the stated prior), with the surviving ray of larger
  couplings closed on each (observationally on the face-value branch, by the
  naturalness prior on Mannheim's).
\item \textbf{Horn 2} (Lorentz $+$ configurations $\Rightarrow$ no
  coupled ghosts): no conformal gravity coupled to matter---the
  theory abandons its purpose.
\item \textbf{Horn 3} (coupled ghosts $+$ Lorentz $\Rightarrow$ no
  configurations): the Bender--Mannheim quantization. This is where
  the theory already lives, by kinematics rather than choice; its
  price, in this currency, is that the ``unitarity'' retained is the
  conditional bookkeeping of an open process
  (Sec.~\ref{sec:lemma})---at the free/kinematic level, within the stated
  dictionary classes, there is no
  underlying \emph{closed} sample space of definite configurations (the lemma and the tower theorem
  are free-field statements; Sec.~\ref{sec:discussion} scopes them).
\end{itemize}
\end{claim}

\begin{table}[t]
\centering\small
\begin{tabular}{@{}p{0.34\linewidth}cp{0.44\linewidth}@{}}
\toprule
escape route & horn & price, in the process currency\\
\midrule
Lorentz-violating cutoff / preferred frame
\cite{CJM2004,Konnig2016} & 1 &
the $\ag$ fork: vacuous-but-observationally-dead, or a
macroscopic-wavelength cutoff (larger-coupling ray closed by fit or stated prior)\\
Lorentz-violating gravitons \cite{Deruelle2012} & 1 &
observable blow-up of the primordial GW spectrum (their result; a two-horn special case)\\
decouple the ghost sector & 2 &
no conformal gravity coupled to matter\\
Maldacena boundary conditions \cite{Maldacena2011} & 2
(attempted) & argued insufficient on Minkowski
\cite{HellLustZoupanos2023}\\
$\PT$/zero-norm quantization
\cite{BenderMannheim2008,Mannheim2020} & 3 &
no \emph{closed} sample space within the stated dictionary classes (Lemma~\ref{lem:jordan}); retained unitarity is
conditioned, open-system bookkeeping\\
negative-norm configuration space \cite{SalvioStrumia2016} & 3 &
configurations without a probability measure\\
fakeons \cite{Anselmi2018} & 3 &
definition-level exit; microcausality violation ``survives the
classical limit''\\
Merlin modes \cite{DonoghueMenezes} & 3-adjacent &
sample space kept, causal order paid (a loop-level width construction, outside the tree-level trilemma)\\
metastability / benign ghosts
\cite{Smilga2005,GSTZ2021,SalvioMeta,ErrastiDiez2024,Deffayet2026,EwasiukProfumo2026} &
(premise) & disputes the rate premise rather than picking a horn;
the open QFT rate is conceded ``naively infinite''; the 2026
quantum-mechanical entries sharpen the dispute (unitary evolution
with a well-defined vacuum; exact moment bounds) without touching
the field-theoretic rate\\
\bottomrule
\end{tabular}
\caption{The adjudicated published escape routes of quadratic/conformal gravity
(one cited work concerns massive gravity's Boulware--Deser ghost
\cite{Konnig2016}), assigned to the trilemma's horns. No route
occupies a fourth position that we can identify; the closest
candidate is the Merlin route, which keeps the sample space and
trades causal order instead---a different door out of the same
room, priced in a currency the trilemma does not list but
operational physics does.}
\label{tab:horns}
\end{table}

The \emph{adjudicated} escape routes of quadratic gravity map onto the horns
(two rows only horn-adjacent or premise-level, as flagged in Table~\ref{tab:horns})
with no further remainder we have been able to identify (the recent (2024--2026) entries
flagged in the introduction, the older alternative quantizations \cite{Mostafazadeh2011,IlhanKovner2014}, and the Bohmian route \cite{DeyFring2026}, remain pending;
Table~\ref{tab:horns}): Lorentz-violating cutoffs and frames
\cite{CJM2004,Konnig2016} (horn 1); Maldacena's boundary conditions
\cite{Maldacena2011}, argued insufficient on Minkowski
\cite{HellLustZoupanos2023} (an attempted horn-2 residence);
$\PT$/zero-norm quantization \cite{BenderMannheim2008,Mannheim2020},
fakeons \cite{Anselmi2018}, and negative-norm configuration space
\cite{SalvioStrumia2016} (horn 3, in three dialects: all exit the
sample space, paying respectively in nonstationary-state bookkeeping,
microcausality that ``survives the classical limit,'' and
probability-free configurations); Merlin modes
\cite{DonoghueMenezes} (horn 3 adjacent: the sample space is kept
but causal order is not, which in the operational terms of
\cite{instantaneous} is a different door out of the same room);
and the trade of Weyl ghosts for observable Lorentz violation in the
primordial gravitational-wave spectrum \cite{Deruelle2012} is a
horn-1 special case (with a horn-2 aspect) computed before the general pattern was
visible. The closest prior statements of a forced choice---%
Donoghue--Menezes' unitarity at the price of microcausality
(violated on scales of order the resonance
width) \cite{DonoghueMenezes} and CJM's stability-versus-Lorentz
\cite{CJM2004}---each name one edge, and Steinwachs names a
different triangle altogether
(unitarity--renormalizability--relativistic
invariance)~\cite{Steinwachs2020}; the present triangle, and the
identification of its third vertex with the existence of a
configuration sample space, appears to be new here, as does every
use of the process currency in this corpus.

%=====================================================================
\section{Discussion: what is established, at what strength}
\label{sec:discussion}

\emph{Theorem-grade}: Lemma~\ref{lem:jordan} and its corollary that
the conformal point's Jordan (no-ghost) dynamics admits no closed
beable representation (per
sector; on the Fock tower for dictionaries boundedly invertible in
the occupation-number Fock norm);
the sector-tower proposition and the bounded-dictionary theorem on
the free Fock space (Sec.~\ref{sec:tower}); the
validity of the open (dilated) process for the mechanical Jordan
cell (Sec.~\ref{sec:lemma}).
\emph{Measured}: the reduced-measure slopes (validated against
\cite{instantaneous}); the per-mode dictionary cost
\eqref{eq:condscaling};
the UV approach to the exceptional line with its units-honest
decomposition; the conditional-mass-creation indivisibility
mechanism of the dipole cell (under the minimal $\sigma_{\max}$
convention, Sec.~\ref{sec:lemma}).
\emph{Argued}: the adjudication of Sec.~\ref{sec:lemma} (that
Bender--Mannheim's nonstationary-state unitarity \emph{is} the
dilation's conditional dynamics is an identification we find exact
in structure, but it is an interpretation of their formalism, not a
theorem about it); the one-loop pole statement of
Sec.~\ref{sec:kinematics} that the Merlin width does not detune
the theory off the exceptional line (a resonance-type EP with
unequal linewidths---an analytic one-loop estimate); the $T\sim s$
vertex parallel behind the contact
rung, and with it the reduced-measure power counting (its
$s$-flatness is forced by scale invariance and its slopes are
measured, but the vertex class itself is the argued premise); the
horn assignments of Sec.~\ref{sec:trilemma}.
\emph{Conditional}: everything in Sec.~\ref{sec:rate}, on the stated
branch structure---in particular the $|\ag|\sim3\times10^{109}$
face-value matching is our arithmetic on published fit parameters,
not a published number, and Mannheim disputes the matching's
premise; we present the fork, not a verdict on it.

Limits. All rate statements are tree-level golden rule on flat
space with Dirac-norm quantization; Mannheim contends flat space
with those norms is the wrong arena twice over, and our reply is the
structure of the paper: on his arena the trilemma's third seat is
already gone (Sec.~\ref{sec:lemma}), so the rate section prices only
the alternative he rejects. The energy-density bound is
deliberately crude (its calibration against the phantom case is
order-of-magnitude); a dedicated soft-photon analysis would tighten
but not rescue horn 1. The equal-frequency quantum mechanics we
adjudicate is no longer only the mechanical prototype---the sector
tower (Sec.~\ref{sec:tower}) carries the no-representation
statement to the free Fock space of every conformal mode, for any
dictionary boundedly invertible in the occupation-number Fock
norm---but the field-theoretic zero-norm construction
still has structure (conformal ghosts on curved backgrounds, the
$S(x)$ sector, interactions) that the free flat-space tower does
not see, and our theorems constrain the latter only.
Finally, the trilemma names combinations, not truths: conformal
gravity may yet be the right theory of nature---the claim is that
if so, it is a theory without a configuration sample space, and its
advertised unitarity is of the conditioned, open-system kind.

%=====================================================================
\phantomsection
\section*{Author contributions and use of AI}
\addcontentsline{toc}{section}{Author contributions and use of AI}

This manuscript was produced through a collaboration, directed by
the author, with the large language model Claude Fable~5
(Anthropic), which carried out the technical work: the research
program, the derivations and proofs, the numerical experiments, the
literature analysis, and the text, including this statement. J.T.\
posed the research direction, provided iterative review, and takes
responsibility for the decision to circulate. The research notebook
in the program's public repository records the development,
including a same-day correction of this paper's own initial
cutoff-bound claim (superseded by the $\ag$ fork of
Sec.~\ref{sec:rate}) when the coupling's face value was pinned down.

%=====================================================================
\phantomsection
\section*{Code availability}
\label{sec:code}
\addcontentsline{toc}{section}{Code availability}

All numerical claims are reproducible from five standalone scripts
(NumPy/SciPy, seconds to minutes, controls printed inline) in the
program's main repository, with companion-paper numbers cited as
such, the face-value $\ag$ matching emitted by
\texttt{scripts/t4\_rate\_ladder.py} from the constants documented in
\texttt{research/t4\_vertex\_conventions.md}, and the one-loop pole
estimates of Sec.~\ref{sec:kinematics} (width, suppression factor,
fixed-linewidth foil) analytic computations whose arithmetic is
stated in-text:
\begin{center}
\url{https://github.com/embeddetech/p05-choosing-ghosts}
\end{center}
(also archived, with the full program history, in the hub
repository \url{https://github.com/embeddetech/ghosts-program});
\texttt{scripts/t4\_\allowbreak ep\_\allowbreak dictionary.py} (Sec.~\ref{sec:kinematics}),
\texttt{scripts/t4\_\allowbreak jordan\_\allowbreak process.py} (Sec.~\ref{sec:lemma}),
\texttt{scripts/t4\_\allowbreak sector\_\allowbreak tower.py} and \texttt{scripts/t4\_\allowbreak tower\_\allowbreak similarity.py} (Sec.~\ref{sec:tower}),
\texttt{scripts/t4\_\allowbreak rate\_\allowbreak ladder.py} (Sec.~\ref{sec:rate}); the
supporting literature and convention reports are under
\texttt{research/}.

%=====================================================================
\begin{thebibliography}{99}\small

\bibitem{indivisible}
J.~Torkelson,
\emph{Indivisible stochastic processes for non-Hermitian and
higher-derivative quantum systems: detailed balance, broken $\PT$
symmetry, and the explosion of ghosts},
companion paper (2026),
\url{https://github.com/embeddetech/p01-indivisible-ghosts}.

\bibitem{instantaneous}
J.~Torkelson,
\emph{Ghost vacuum decay as a stochastic process: instantaneous
states, the boost-orbit measure, and the foliation that ghosts make
visible},
companion paper (2026),
\url{https://github.com/embeddetech/p03-instantaneous-ghosts}.

\bibitem{conditioned}
J.~Torkelson,
\emph{The price of a post-selected signal: a sharp
survival--distinguishability trade-off, saturated by broken-$\PT$
dynamics},
companion paper (2026),
\url{https://github.com/embeddetech/p04-conditioned-ghosts}.

\bibitem{Barandes2023}
J.~A.~Barandes,
\emph{The stochastic-quantum correspondence},
Philos.\ Phys.\ \textbf{3}(1), 8 (2025), arXiv:2302.10778.

\bibitem{Barandes2025}
J.~A.~Barandes,
\emph{Quantum systems as indivisible stochastic processes},
arXiv:2507.21192.

\bibitem{PaisUhlenbeck1950}
A.~Pais and G.~E.~Uhlenbeck,
\emph{On field theories with non-localized action},
Phys.\ Rev.\ \textbf{79}, 145 (1950).

\bibitem{BenderMannheim2008}
C.~M.~Bender and P.~D.~Mannheim,
\emph{No-ghost theorem for the fourth-order derivative
Pais--Uhlenbeck oscillator model},
Phys.\ Rev.\ Lett.\ \textbf{100}, 110402 (2008), arXiv:0706.0207.

\bibitem{BenderMannheim2008PRD}
C.~M.~Bender and P.~D.~Mannheim,
\emph{Exactly solvable PT-symmetric Hamiltonian having no Hermitian
counterpart},
Phys.\ Rev.\ D \textbf{78}, 025022 (2008), arXiv:0804.4190.

\bibitem{Mannheim2011}
P.~D.~Mannheim,
\emph{Making the case for conformal gravity},
Found.\ Phys.\ \textbf{42}, 388 (2012), arXiv:1101.2186.

\bibitem{Mannheim2020}
P.~D.~Mannheim,
\emph{Ghost problems from Pauli--Villars to fourth-order quantum
gravity and their resolution},
Int.\ J.\ Mod.\ Phys.\ D \textbf{29}, 2043009 (2020),
arXiv:2004.00376.

\bibitem{Mannheim2015}
P.~D.~Mannheim,
\emph{Comment on ``Problems with Mannheim's conformal gravity
program''},
Phys.\ Rev.\ D \textbf{93}, 068501 (2016), arXiv:1506.02479.

\bibitem{MannheimOBrien2011}
P.~D.~Mannheim and J.~G.~O'Brien,
\emph{Fitting galactic rotation curves with conformal gravity and a
global quadratic potential},
Phys.\ Rev.\ D \textbf{85}, 124020 (2012), arXiv:1011.3495.

\bibitem{Smilga2005}
A.~V.~Smilga,
\emph{Benign vs.\ malicious ghosts in higher-derivative theories},
Nucl.\ Phys.\ B \textbf{706}, 598 (2005), arXiv:hep-th/0407231.

\bibitem{Smilga2008}
A.~V.~Smilga,
\emph{Comments on the dynamics of the Pais--Uhlenbeck oscillator},
SIGMA \textbf{5}, 017 (2009), arXiv:0808.0139.

\bibitem{Smilga2014}
A.~V.~Smilga,
\emph{Exceptional points of infinite order giving a continuous
spectrum},
Int.\ J.\ Theor.\ Phys.\ \textbf{54}, 3900 (2015), arXiv:1409.8450.

\bibitem{FMT2026}
A.~Fring, I.~Marquette, and T.~Taira,
\emph{Spectrum-generating algebra and intertwiners of the resonant
Pais--Uhlenbeck oscillator},
arXiv:2601.20752 (their abstract: classically equivalent
Hamiltonians define inequivalent quantum theories).

\bibitem{FringTairaTurner2025}
A.~Fring, T.~Taira, and B.~Turner,
\emph{Ghost-free quantisation of higher time-derivative theories
via non-unitary similarity transformations},
arXiv:2506.21400 (2025).

\bibitem{FelskiFringTurner2025}
A.~Felski, A.~Fring, and B.~Turner,
\emph{Lie symmetries and ghost-free representations of the
Pais--Uhlenbeck model},
Mod.\ Phys.\ Lett.\ A \textbf{41}, 2650019 (2026), arXiv:2505.07869.

\bibitem{DeyFring2026}
S.~Dey and A.~Fring,
\emph{Quantum--classical diagnostics and Bohmian inequivalence for
higher time-derivative Hamiltonians},
Universe \textbf{12}(7), 200 (2026), arXiv:2603.15395.

\bibitem{Holdom2024}
B.~Holdom,
\emph{Making sense of ghosts},
Nucl.\ Phys.\ B \textbf{1008}, 116696 (2024), arXiv:2408.04089.

\bibitem{Kuntz2024}
J.~Kuntz,
\emph{Unitarity through PT symmetry in quantum quadratic gravity},
Class.\ Quantum Grav.\ \textbf{42}, 175003 (2025), arXiv:2410.08278.

\bibitem{BatemanTurok2026}
S.~Bateman and N.~Turok,
\emph{Escape from Ostrogradsky via Hidden Ghost Parity},
arXiv:2607.00096.

\bibitem{KumarMarto2026}
K.~S.~Kumar and J.~Marto,
\emph{Quantum (quadratic) gravity: replacing the massive tensor
ghost with an inverted harmonic oscillator-like instability},
arXiv:2603.07150;
\emph{Unitary Quadratic Quantum Gravity in 4D},
arXiv:2604.19707.

\bibitem{AsoreyKPS2026}
M.~Asorey, G.~Krein, M.~Pardina, and I.~L.~Shapiro,
\emph{Reflection positivity in a higher-derivative model with
physical bound states of ghosts},
JHEP \textbf{02} (2026) 020, arXiv:2511.15283;
\emph{Bound states of massive complex ghosts in superrenormalizable
quantum gravity theories},
JHEP \textbf{01} (2025) 113, arXiv:2408.16514.

\bibitem{Santillan2025}
O.~P.~Santill\'an,
\emph{LSZ ghostbusters in the quadratic gravity stage},
Eur.\ Phys.\ J.\ C \textbf{85}, 810 (2025), arXiv:2505.20597.

\bibitem{Buoninfante2026}
L.~Buoninfante,
\emph{Asymptotic Quantum Dynamics of Ghost Fields},
arXiv:2605.29047.

\bibitem{ErrastiDiez2024}
V.~Errasti D\'iez, J.~Gaset Rif\`a, and G.~Staudt,
\emph{Foundations of ghost stability},
Fortschr.\ Phys.\ \textbf{73}, 2400268 (2025), arXiv:2408.16832.

\bibitem{Deffayet2026}
C.~Deffayet, A.~Fathe Jalali, A.~Held, S.~Mukohyama, and A.~Vikman,
\emph{Unitary time evolution and vacuum for a quantum stable ghost},
arXiv:2604.21823.

\bibitem{EwasiukProfumo2026}
C.~Ewasiuk and S.~Profumo,
\emph{Ghost degrees of freedom without quantum runaway: exact moment
bounds from an operator conservation law},
arXiv:2604.21348.

\bibitem{Oda2026}
I.~Oda,
\emph{Confinement of Massive Ghost in Quadratic Gravity},
arXiv:2605.01966.

\bibitem{Mostafazadeh2011}
A.~Mostafazadeh,
\emph{A Hamiltonian formulation of the Pais--Uhlenbeck oscillator
that yields a stable and unitary quantum system},
Phys.\ Lett.\ A \textbf{375}, 93 (2010), arXiv:1008.4678.

\bibitem{GSTZ2021}
C.~Gross, A.~Strumia, D.~Teresi, and M.~Zirilli,
\emph{Is negative kinetic energy metastable?},
Phys.\ Rev.\ D \textbf{103}, 115025 (2021), arXiv:2007.05541.

\bibitem{CJM2004}
J.~M.~Cline, S.~Jeon, and G.~D.~Moore,
\emph{The phantom menaced: constraints on low-energy effective
ghosts},
Phys.\ Rev.\ D \textbf{70}, 043543 (2004), hep-ph/0311312.

\bibitem{ClinePhantom2024}
J.~M.~Cline, M.~Puel, T.~Toma, and Q.~S.~Wang,
\emph{Phantom fluid cosmology: impact of a phantom hidden sector on
cosmological observables},
Phys.\ Rev.\ D \textbf{108}, 095033 (2023), arXiv:2308.12989.

\bibitem{SalvioStrumia2016}
A.~Salvio and A.~Strumia,
\emph{Quantum mechanics of 4-derivative theories},
Eur.\ Phys.\ J.\ C \textbf{76}, 227 (2016), arXiv:1512.01237.

\bibitem{Salvio2018}
A.~Salvio,
\emph{Quadratic gravity},
Front.\ Phys.\ \textbf{6}, 77 (2018), arXiv:1804.09944.

\bibitem{SalvioMeta}
A.~Salvio,
\emph{Metastability in quadratic gravity},
Phys.\ Rev.\ D \textbf{99}, 103507 (2019), arXiv:1902.09557.

\bibitem{Mannheim2007}
P.~D.~Mannheim,
\emph{Schwarzschild limit of conformal gravity in the presence of
macroscopic scalar fields},
Phys.\ Rev.\ D \textbf{75}, 124006 (2007), arXiv:gr-qc/0703037.

\bibitem{DonoghueMenezes}
J.~F.~Donoghue and G.~Menezes,
\emph{Unitarity, stability, and loops of unstable ghosts},
Phys.\ Rev.\ D \textbf{100}, 105006 (2019), arXiv:1908.02416;
\emph{The arrow of causality and quantum gravity},
Phys.\ Rev.\ Lett.\ \textbf{123}, 171601 (2019), arXiv:1908.04170.

\bibitem{Anselmi2018}
D.~Anselmi,
\emph{On the quantum field theory of the gravitational
interactions},
JHEP \textbf{06} (2017) 086, arXiv:1704.07728;
\emph{Fakeons and Lee--Wick models},
JHEP \textbf{02} (2018) 141, arXiv:1801.00915;
\emph{Fakeons, microcausality and the classical limit of quantum
gravity},
Class.\ Quantum Grav.\ \textbf{36}, 065010 (2019), arXiv:1809.05037.

\bibitem{HellLustZoupanos2023}
A.~Hell, D.~L\"ust, and G.~Zoupanos,
\emph{On the ghost problem of conformal gravity},
JHEP \textbf{08} (2023) 168, arXiv:2306.13714.

\bibitem{IlhanKovner2014}
I.~B.~Ilhan and A.~Kovner,
\emph{Some comments on ghosts and unitarity: the Pais--Uhlenbeck
oscillator revisited},
Phys.\ Rev.\ D \textbf{88}, 044045 (2013), arXiv:1301.4879
[see also \emph{Conformal gravity redux: Ghost-turned tachyon},
Phys.\ Rev.\ D \textbf{89}, 104015 (2014), arXiv:1403.4664].

\bibitem{Flanagan2006}
\'E.~\'E.~Flanagan,
\emph{Fourth order Weyl gravity},
Phys.\ Rev.\ D \textbf{74}, 023002 (2006), astro-ph/0605504.

\bibitem{Yoon2013}
Y.~Yoon,
\emph{Problems with Mannheim's conformal gravity program},
Phys.\ Rev.\ D \textbf{88}, 027504 (2013), arXiv:1305.0163.

\bibitem{Riegert1984}
R.~J.~Riegert,
\emph{The particle content of linearized conformal gravity},
Phys.\ Lett.\ A \textbf{105}, 110 (1984).

\bibitem{Maldacena2011}
J.~Maldacena,
\emph{Einstein gravity from conformal gravity},
arXiv:1105.5632.

\bibitem{Deruelle2012}
N.~Deruelle, M.~Sasaki, Y.~Sendouda, and A.~Youssef,
\emph{Lorentz-violating vs.\ ghost gravitons: the example of Weyl
gravity},
JHEP \textbf{09} (2012) 009, arXiv:1202.3131.

\bibitem{Steinwachs2020}
C.~F.~Steinwachs,
\emph{Towards a unitary, renormalizable, and ultraviolet-complete
quantum theory of gravity},
Front.\ Phys.\ \textbf{8}, 185 (2020), arXiv:2004.07842.

\bibitem{Konnig2016}
F.~K\"onnig, H.~Nersisyan, Y.~Akrami, L.~Amendola, and
M.~Zumalac\'arregui,
\emph{A spectre is haunting the cosmos: quantum stability of massive
gravity with ghosts},
JHEP \textbf{11} (2016) 118, arXiv:1605.08757.

\end{thebibliography}

\end{document}
